\providecommand{\Res}[2]{\qty[round-mode=figures,round-precision=3]{\fpeval{#1}}{#2}}
\providecommand{\Val}[1]{\num[round-mode=figures,round-precision=4]{\fpeval{#1}}}
\providecommand{\ResSci}[2]{\qty[round-mode=figures,round-precision=3,exponent-mode=scientific]{\fpeval{#1}}{#2}}
\subsubsection*{Question a) Force électrique}
\begin{enumerate}[label=\arabic*)]
    \item \textbf{Variables connues et recherchées :} $|q_1| = \qty{\v{q1}e-6}{C}$, $|q_2| = \qty{\v{q2}e-6}{C}$, $r = \qty{\vsi{r}}{m}$, $k = \qty{8.99e9}{N.m^2/C^2}$ ; variable recherchée : $F$
    \item \textbf{Équations pertinentes :} loi de Coulomb :
    \[ F = \frac{k\,|q_1|\,|q_2|}{r^2} \]
    \item \textbf{Isolation de la variable recherchée :} l'équation est déjà isolée.
    \item \textbf{Résolution numérique :}
    \[ F = \frac{\num{8.99e9} \cdot \num{\v{q1}e-6} \cdot \num{\v{q2}e-6}}{(\num{\vsi{r}})^2} = \Res{8.99e9*\vsi{q1}*\vsi{q2}/\vsi{r}^2}{N} \]
    \item \textbf{Réponse :} $F = \mathbf{\Res{8.99e9*\vsi{q1}*\vsi{q2}/\vsi{r}^2}{N}}$.
\end{enumerate}

\subsubsection*{Question b) Nature de la force}
Les charges sont de signes contraires (une positive, une négative) : la force est \textbf{attractive}.

\subsubsection*{Question c) Distance pour une force de \qty{\v{Fc}}{N}}
\begin{enumerate}[label=\arabic*)]
    \item \textbf{Variables connues et recherchées :} $F = \qty{\v{Fc}}{N}$, $|q_1|$, $|q_2|$, $k$ ; variable recherchée : $r$
    \item \textbf{Équations pertinentes :}
    \[ F = \frac{k\,|q_1|\,|q_2|}{r^2} \]
    \item \textbf{Isolation de la variable recherchée :}
    \[ r = \sqrt{\frac{k\,|q_1|\,|q_2|}{F}} \]
    \item \textbf{Résolution numérique :}
    \[ r = \sqrt{\frac{\num{8.99e9} \cdot \num{\v{q1}e-6} \cdot \num{\v{q2}e-6}}{\num{\v{Fc}}}} = \Res{sqrt(8.99e9*\vsi{q1}*\vsi{q2}/\v{Fc})}{m} \]
    \item \textbf{Réponse :} il faut placer les charges à \textbf{\Res{100*sqrt(8.99e9*\vsi{q1}*\vsi{q2}/\v{Fc})}{cm}} l'une de l'autre.
\end{enumerate}
